\section{Introduction}
\label{sec:introduction}

PHP continues to power a large portion of production web software, much of it in long-lived legacy systems.\cite{w3techs_php,w3techs_wordpress,wordpress}
Migrating these systems to newer PHP versions is difficult.
Required changes are scattered across the codebase and span evolving language features, library updates, and stricter typing constraints.\cite{php83spec,persson2015php}
Developers need visibility not only into whether a migration succeeds, but also into which upgrade requirements are satisfied and which remain unresolved.

Large language models are increasingly used to assist software modernization, yet evaluating their effectiveness remains challenging.
Most existing studies rely on outcome signals such as parsing success, linting results, or test outcomes.\cite{openai2023gpt4,roziere2023codellama,cassano2024canitedit}
These signals are useful, but they provide limited insight into what changed during migration.
They show whether code runs, but they do not reveal which upgrade obligations were resolved or newly introduced.
This limits systematic comparison across models and prompting strategies, especially when performance varies with context limits and input size.\cite{liu2024lost,du2025context,jimenez2024swebench}

We introduce \textbf{Pinned Contract Evaluation (PCE)}, a policy-grounded framework for evaluating automated code modernization.
PCE treats a static analysis upgrade pipeline as a versioned migration contract.
Executing the contract on the original program induces an obligation space defined by triggered upgrade rule types.
Running the same contract on migrated outputs reveals which obligations are discharged, which remain, and which new ones appear.
Because the policy is fixed, results are directly comparable across systems and experiments.
PCE evaluates progress relative to a fixed upgrade policy rather than attempting to establish full behavioral or semantic correctness.
Unlike prior evaluations that treat static-analysis outputs only as coarse outcome signals, PCE formalizes the upgrade policy itself as a contract that induces a measurable obligation space for systematic comparison across systems.

To contextualize contract-level progress, we include additional diagnostics.
These include syntax validation, structural preservation signals, compatibility analysis using PHPCompatibility, and a regression-oriented loadability sentinel.

We instantiate PCE for PHP version migration using a pinned Rector upgrade pipeline.\cite{rector}
Across five LLM systems evaluated on a 100-file WordPress 4.0 benchmark, mean obligation discharge ranges from 38\% to 78\%.
Remaining work concentrates in PHP 8 rule families, and high discharge does not necessarily imply syntactic or runtime robustness.
These findings show why policy-grounded evaluation is useful for understanding both progress and failure modes in automated modernization.

\paragraph{Contributions}
This paper makes three contributions.

\begin{enumerate}[leftmargin=*]

\item \textbf{Pinned Contract Evaluation (PCE).}
A reproducible evaluation framework that models modernization as progress through a contract-induced obligation space defined by a pinned static analysis policy.

\item \textbf{Contract-induced migration benchmark.}
A 100-file WordPress 4.0 dataset that preserves the obligation space induced by the contract.

\item \textbf{Reproducible evaluation harness.}
A standardized LLM migration and scoring pipeline with deterministic extraction, reconstruction, and fully logged artifacts.

\end{enumerate}

\subsection{Research Questions}
\label{sec:rqs}

Under a fixed contract and its induced obligation space, we ask:

\begin{enumerate}[leftmargin=*]

\item RQ1. How do models compare in obligation discharge and residual obligation profiles?

\item RQ2. Which rules and language milestones remain most resistant to automated migration?

\item RQ3. How does contract progress relate to syntax validity, structural preservation, and loadability regressions?

\item RQ4. How do PCE outcomes align with results from an independent compatibility analysis using PHPCompatibility?

\end{enumerate}

\subsection{Artifacts and Code Availability}
\label{subsec:artifacts}

We release all prompts, configurations, benchmark metadata, and evaluation scripts in a public repository to enable independent reproduction of our results.\cite{wohlin2012experimentation}